\chapter{Conclusões} \label{chap:conclusoes}

% Objetivo: sintetizar achados, contribuições, limitações e agenda futura.

\section{Síntese dos achados}
% TODO: aguardando a rodada confirmatória (vereditos H1/H2 por horizonte e
% contribuição de famílias H3). Preencher após consolidação dos artefatos.

\section{Contribuições do trabalho}
As contribuições deste trabalho organizam-se em duas frentes: metodológicas, já consolidadas pelo desenho experimental, e empíricas, dependentes da execução da rodada confirmatória.

\subsection{Contribuições metodológicas}
A primeira contribuição é o enquadramento da avaliação de previsões de retornos diários na distribuição preditiva, e não no ponto previsto: a calibração é o objeto primário e, ao mesmo tempo, \textit{gate} de elegibilidade para a comparação de \textit{skill}, e a refutação das hipóteses é tratada como resultado informativo.

A segunda é o tratamento da saída quantílica de regressores treinados por nível. Todos os modelos emitem a mesma grade densa de quantis; o rearranjo monotônico, fundamentado no resultado de \cite{CHERNOZHUKOV2010REARRANGEMENT}, é aplicado como parte do estimador, com os valores bruto e rearranjado persistidos por nível para que seu efeito seja auditável; e a degeneração da grade é tratada por um \textit{gate} separado, de modo que o rearranjo não mascare previsões sem informação distribucional.

A terceira é uma hierarquia de \textit{baselines} em três níveis --- naive, estatísticos e \textit{gradient boosting} quantílico ---, todos emitindo a mesma grade sob o mesmo protocolo temporal e o mesmo protocolo de informação do candidato, com convenções de emissão de quantis declaradas antecipadamente. A presença de um comparador forte de aprendizado de máquina eleva a barra de H2 e torna o empate um resultado honesto.

A quarta é o protocolo temporal: validação \textit{walk-forward} de janela expansiva com purga e embargo contados em sessões de pregão e um conjunto de calibração dedicado, nunca usado para ajuste ou seleção, que viabiliza a predição conformal como referência de cobertura empírica.

A quinta é a disciplina confirmatória: candidato único \textit{all-features} sem seleção por desempenho fora da amostra, pré-registro imutável \textit{hasheado} e um \textit{scorecard} aplicado mecanicamente, que separa o veredito do perfil comparativo e nunca agrega resultados entre horizontes.

A sexta é a hierarquia de evidência para a contribuição de famílias de indicadores --- VSN, importância por permutação e ablação, com consistência em pelo menos dois de três métodos --- e a distinção formal entre contribuição preditiva condicional ao modelo e inferência causal.

Por fim, a implementação, descrita no Apêndice D, valida cada métrica e teste estatístico contra implementações de referência e permite reconstruir todas as conclusões a partir das previsões persistidas, sem re-treino.

\subsection{Contribuições empíricas}
% TODO: consolidar contribuições empíricas (calibração observada por horizonte;
% comparação TFT vs baselines; contribuição relativa de famílias).
As contribuições empíricas reportarão, sob o protocolo acima, a calibração probabilística observada por horizonte, a comparação pareada do candidato contra a hierarquia de \textit{baselines} e a contribuição relativa das famílias de indicadores. Sua consolidação depende da execução da rodada confirmatória, sob o pré-registro, e será incorporada nesta seção em seguida.

\section{Limitações}
As limitações deste trabalho decorrem, em parte, do recorte assumido e devem ser interpretadas conjuntamente com os achados, independentemente do desfecho das hipóteses.

Em primeiro lugar, os modelos são \textit{single-asset}: cada execução é treinada e avaliada para um ativo específico, com AAPL como piloto. O procedimento é replicável para outros ativos, mas as conclusões valem apenas para o ativo e o período avaliados.

Em segundo lugar, as análises de contribuição relativa de famílias indicam contribuição preditiva condicional ao modelo treinado, e não causalidade. Interpretações locais são apenas ilustrativas; conclusões derivam de evidência global por horizonte e da estabilidade entre sementes e \textit{folds}.

Em terceiro lugar, retornos diários de ativos líquidos têm previsibilidade muito baixa na média \cite{GU2020}, o que torna plausível a ausência de superioridade do candidato frente a comparadores fortes. A não-dominância em H2 é, nesse contexto, resultado científico válido sobre os limites do estimador, e não falha do método.

Em quarto lugar, o horizonte \(h+30\) é tratado como suplementar, dada a amostra efetiva menor e a sobreposição temporal entre observações; as conclusões confirmatórias concentram-se em \(h+1\) e \(h+7\).

Em quinto lugar, os resultados da rodada exploratória de hiperparâmetros não constituem evidência; as conclusões dependem exclusivamente da rodada confirmatória com hiperparâmetros congelados \cite{RASCHKA2018}.

Em sexto lugar, a garantia de cobertura da predição conformal pressupõe permutabilidade, que não vale em séries temporais; a cobertura dos intervalos conformais é, por isso, reportada como empírica \cite{BARBER2023}.

Por fim, as bandas de calibração e os limiares de qualidade não têm valor canônico na literatura; sua validade decorre da declaração antecipada no pré-registro, conforme discutido na Seção~\ref{sec:scorecard}.

\section{Trabalhos futuros}
Como evoluções coerentes com as limitações observadas, propõem-se cinco direções.

A primeira é ampliar a avaliação para outros ativos, mantendo o desenho \textit{single-asset} e definindo no pré-registro os critérios de inclusão de ativos, com declaração explícita de disponibilidade, qualidade e comparabilidade.

A segunda é tratar a classificação ou detecção de regimes de mercado como estudo separado, com rotulagem e metodologia próprias; o presente trabalho registra regimes de alta volatilidade apenas como ameaça à validade, não como objeto de modelagem.

A terceira é explorar variantes de emissão de quantis e de calibração não adotadas aqui, como caudas t-Student para o EWMA, quantis não cruzantes por construção \cite{BONDELL2010,CANNON2018} e variantes conformais adaptativas para séries temporais.

A quarta é o desenho de estudo específico para inferência causal das famílias de indicadores; o escopo atual sustenta apenas contribuição preditiva condicional ao modelo.

A quinta é a aplicação do estimador em alocação de carteira ou estratégia de \textit{trading}, e a extensão a outras classes de ativos e frequências (por exemplo, derivativos de criptoativos e dados de microestrutura), condicionadas a validação própria --- custos de transação, restrições de execução --- fora do recorte deste trabalho.
